% texte.tex — f164v, vv. 165–207 (conversion initiale des travaux antérieurs)
\vers[n=165]{}{}{αὐτός τε καλὸς ἦν καὶ καλῶς ἠμπέσχετο}{beau lui-même, il s'habillait de beaux vêtements}
\vers[n=166]{}{}{διὰ τοῦτ᾽ ἄρ᾽ αὐτοῦ καὶ κάλ᾽ ἦν τὰ δράματα.}{et c'est aussi ce qui fait que ses drames étaient beaux.}
\vers[n=167]{}{}{ὅμοια γὰρ ποιεῖν ἀνάγκη τῇ φύσει.}{Car c'est une nécessité que l'on compose suivant sa nature.}
\vers[n=168]{ΜΝ.}{Mnésiloque}{ταῦτ᾽ ἄρ᾽ ὁ Φιλοκλέης αἰσχρὸς ὢν αἰσχρῶς ποιεῖ,}{Voilà aussi pourquoi Philoclès, qui est laid, compose de laides choses ;}
\vers[n=169]{}{}{ὁ δ᾽ αὖ Ξενοκλέης ὢν κακὸς κακῶς ποιεῖ,}{pourquoi Xénoclès, qui est mauvais, en compose de mauvaises ;}
\vers[n=170]{}{}{ὁ δ᾽ αὖ Θέογνις ψυχρὸς ὢν ψυχρῶς ποιεῖ.}{pourquoi Théognis, qui est froid, compose froidement.}
\vers[n=171]{ΑΓ.}{Agathon}{ἅπασ᾽ ἀνάγκη ταῦτα γάρ τοι γνοὺς ἐγὼ}{De toute rigueur. C'est, en effet, pour avoir su cela}
\vers[n=172]{}{}{ἐμαυτὸν ἐθεράπευσα.}{que j'ai soigné ma personne.}
\vers[suite]{ΜΝ.}{Mnésiloque}{πῶς πρὸς τῶν θεῶν;}{Comment, au nom des dieux ?}
\vers[n=173]{ΕΥ.}{Euripide}{παῦσαι βαΰζων καὶ γὰρ ἐγὼ τοιοῦτος ἦν}{Cesse de japper. Car tel j'étais aussi moi-même}
\vers[n=174]{}{}{ὢν τηλικοῦτος, ἡνίκ᾽ ἠρχόμην ποιεῖν.}{à son âge, quand je commençais à produire.}
\vers[n=175]{ΜΝ.}{Mnésiloque}{μὰ τὸν Δί᾽ οὐ ζηλῶ σε τῆς παιδεύσεως.}{Par Zeus ! je ne t'envie pas ton éducation.}
\vers[n=176]{ΕΥ.}{Euripide}{ἀλλ᾽ ὧνπερ οὕνεκ᾽ ἦλθον, ἔα μ᾽ εἰπεῖν.}{Mais laisse-moi dire ce qui m'amène.}
\vers[suite]{ΑΓ.}{Agathon}{λέγε.}{Parle.}
\vers[n=177]{ΕΥ.}{Euripide}{Ἀγάθων, σοφοῦ πρὸς ἀνδρός, ὅστις ἐν βραχεῖ}{Agathon, « c'est le propre d'un sage de savoir résumer en bref}
\vers[n=178]{}{}{πολλοὺς καλῶς οἷός τε συντέμνειν λόγους.}{de longs discours ».}
\vers[n=179]{}{}{ἐγὼ δὲ καινῇ ξυμφορᾷ πεπληγμένος}{Frappé d'une calamité inouïe,}
\vers[n=180]{}{}{ἱκέτης ἀφῖγμαι πρὸς σέ.}{je viens à toi en suppliant.}
\vers[suite]{ΑΓ.}{Agathon}{τοῦ χρείαν ἔχων;}{Que te faut-il ?}
\vers[n=181]{ΕΥ.}{Euripide}{μέλλουσί μ᾽ αἱ γυναῖκες ἀπολεῖν τήμερον}{Les femmes pensent consommer ma perte, aujourd'hui,}
\vers[n=182]{}{}{τοῖς Θεσμοφορίοις, ὅτι κακῶς αὐτὰς λέγω.}{aux Thesmophories, parce que je dis du mal d'elles.}
\vers[n=183]{ΑΓ.}{Agathon}{τίς οὖν παρ᾽ ἡμῶν ἐστιν ὠφέλειά σοι;}{Et que pouvons-nous pour toi ?}
\vers[n=184]{ΕΥ.}{Euripide}{ἡ πᾶσ᾽· ἐὰν γὰρ ἐγκαθεζόμενος λάθρᾳ}{Tout. Si siégeant en tapinois parmi les femmes,}
\vers[n=185]{}{}{ἐν ταῖς γυναιξίν, ὡς δοκῶν εἶναι γυνή,}{en feignant d'être l'une d'elles,}
\vers[n=186]{}{}{ὑπεραποκρίνῃ μου, σαφῶς σώσεις ἐμέ.}{tu prends ma défense, tu ne manqueras pas à me sauver.}
\vers[n=187]{}{}{μόνος γὰρ ἂν λέξειας ἀξίως ἐμοῦ.}{Car seul tu parleras d'une manière digne de moi.}
\vers[n=188]{ΑΓ.}{Agathon}{ἔπειτα πῶς οὐκ αὐτὸς ἀπολογεῖ παρών;}{Alors, que ne vas-tu te défendre en personne ?}
\vers[n=189]{ΕΥ.}{Euripide}{ἐγὼ φράσω σοι. πρῶτα μὲν γιγνώσκομαι}{Je te le dirai. D'abord on me connaît.}
\vers[n=190]{}{}{ἔπειτα πολιός εἰμι καὶ πώγων᾽ ἔχω,}{Ensuite je suis tout chenu, et j'ai de la barbe.}
\vers[n=191]{}{}{σὺ δ᾽ εὐπρόσωπος λευκὸς ἐξυρημένος}{Toi, tu es avenant de figure, blanc, rasé,}
\vers[n=192]{}{}{γυναικόφωνος ἁπαλὸς εὐπρεπὴς ἰδεῖν.}{doué d'une voix de femme, mignon, plaisant à voir.}
\vers[n=193]{ΑΓ.}{Agathon}{Εὐριπίδη—}{Euripide...}
\vers[suite]{ΕΥ.}{Euripide}{τί ἔστιν;}{Eh bien ?}
\vers[suite]{ΑΓ.}{Agathon}{ἐποίησάς ποτε,}{Tu écrivis un jour : « Tu prends plaisir}
\vers[n=194]{}{}{`χαίρεις ὁρῶν φῶς, πατέρα δ᾽ οὐ χαίρειν δοκεῖς;'}{à voir cette lumière, penses-tu que ton père y prenne moins plaisir ? »}
\vers[n=195]{ΕΥ.}{Euripide}{ἔγωγε.}{Sans doute.}
\vers[suite]{ΑΓ.}{Agathon}{μή νυν ἐλπίσῃς τὸ σὸν κακὸν}{N'espère donc pas que je prenne sur mon compte ton infortune.}
\vers[n=196]{}{}{ἡμᾶς ὑφέξειν. καὶ γὰρ ἂν μαινοίμεθ᾽ ἂν.}{Il faudrait que je fusse fou.}
\vers[n=197]{}{}{ἀλλ᾽ αὐτὸς ὅ γε σόν ἐστιν οἰκείως φέρε.}{Charge-toi personnellement de tes propres affaires.}
\vers[n=198]{}{}{τὰς συμφορὰς γὰρ οὐχὶ τοῖς τεχνάσμασιν}{Car les revers, il convient de les subir à force, non d'artifices,}
\vers[n=199]{}{}{φέρειν δίκαιον ἀλλὰ τοῖς παθήμασιν.}{mais d'endurance.}
\vers[n=200]{ΜΝ.}{Mnésiloque}{καὶ μὴν σύ γ᾽ ὦ κατάπυγον εὐρύπρωκτος εἶ}{Au fait toi-même, bougre de cul, tu as l'anus béant à force}
\vers[n=201]{}{}{οὐ τοῖς λόγοισιν ἀλλὰ τοῖς παθήμασιν.}{non de discours mais d'endurance.}
\vers[n=202]{ΕΥ.}{Euripide}{τί δ᾽ ἔστιν ὅτι δέδοικας ἐλθεῖν αὐτόσε;}{Mais qu'est-ce qui te fait craindre d'aller là ?}
\vers[n=203]{ΑΓ.}{Agathon}{κάκιον ἀπολοίμην ἂν ἢ σύ.}{Je finirais plus misérablement que toi.}
\vers[suite]{ΕΥ.}{Euripide}{πῶς;}{Pourquoi ?}
\vers[suite]{ΑΓ.}{Agathon}{ὅπως;}{Pourquoi ?}
\vers[n=204]{}{}{δοκῶν γυναικῶν ἔργα νυκτερείσια}{Je passe pour ravir les travaux nocturnes des femmes,}
\vers[n=205]{}{}{κλέπτειν ὑφαρπάζειν τε θήλειαν Κύπριν.}{et dérober la Cypris féminine.}
\vers[n=206]{ΜΝ.}{Mnésiloque}{ἰδού γε κλέπτειν· νὴ Δία βινεῖσθαι μὲν οὖν.}{Mon cul, ravir ! Eh par Zeus! dis de te faire enfiler.}
\vers[n=207]{}{}{ἀτὰρ ἡ πρόφασίς γε νὴ Δί᾽ εἰκότως ἔχει.}{Mais, par Zeus, la défaite est bonne.}
